\documentclass[11pt]{article}

\usepackage[letterpaper,margin=1in]{geometry}
\usepackage{booktabs}
\usepackage{graphicx}
\usepackage{hyperref}
\usepackage[numbers,sort&compress]{natbib}
\usepackage{microtype}

\hypersetup{
  colorlinks=true,
  linkcolor=black,
  citecolor=black,
  urlcolor=black
}

\newcommand{\draftnote}[1]{\textbf{[Draft note: #1]}}
\newcommand{\placeholderfigure}[2]{%
  \begin{figure}[t]
    \centering
    \fbox{\parbox{0.92\linewidth}{\vspace{0.35in}\centering #2\vspace{0.35in}}}
    \caption{#1}
  \end{figure}
}

\title{GeneLab Benchmark: Cross-Mission Evaluation of Foundation Models and Classical Baselines for Spaceflight Transcriptomics}

\author{
  Jaeyoung Kang\\
  \textit{TODO: affiliation}\\
  \texttt{TODO: email}
}

\date{}

\begin{document}
\maketitle

\begin{abstract}
Spaceflight transcriptomics is an extreme small-sample domain-shift setting:
models must generalize across missions, hardware, duration, and tissue
contexts. We introduce GeneLab Benchmark, a mission-held-out evaluation suite
for NASA Open Science Data Repository mouse spaceflight RNA-seq, and compare
classical baselines, gene-expression foundation models, and text LLMs. Across
the v7.1 public benchmark surface, tuned classical methods remain the strongest
baselines: PCA-LR has the best 8-tissue gene-level mean AUROC, while
gene-expression foundation models and text LLMs fail to consistently outperform
the classical benchmark surface. Pathway-level features reduce mission and
batch separability and help explain tissue-specific transfer, with thymus
generalizing more strongly than liver under cross-mission evaluation. The
benchmark provides reproducible tasks, processed feature matrices, and result
artifacts for evaluating future biological foundation models under real
distribution shift.
\end{abstract}

\section{Introduction}

Space biology has rapidly accumulated public transcriptomic data, but the
field still lacks standardized benchmarks for whether learned spaceflight
signatures generalize across missions. Random sample splits are especially
misleading in this setting because mission identity can couple biological
state, hardware, duration, strain, preprocessing, and cohort structure. A useful
benchmark must therefore test whether a model trained on one set of missions
can classify flight versus ground control in missions it has not seen.

This question is also timely for biological foundation models. Models trained
on large single-cell corpora are often expected to transfer broadly, yet
small-sample bulk RNA-seq under mission shift has different statistical
failure modes than same-distribution cell-level benchmarks. GeneLab Benchmark
tests this setting directly using NASA Open Science Data Repository data
\citep{nasa_osdr}, processed feature matrices, and standardized evaluation
code \citep{genelab_benchmark_repo,genelab_benchmark_hf}.

We make four contributions:

\begin{enumerate}
  \item We define a mission-held-out benchmark for mouse spaceflight
  transcriptomics, with fold construction and feature selection designed to
  reduce cross-mission leakage.
  \item We compare classical machine-learning methods across the v4
  multi-method surface: 8 tissues, 8 classifiers, and 4 feature types.
  \item We summarize gene-expression foundation model and text LLM performance
  against the classical benchmark surface.
  \item We connect prediction performance to pathway-level biology, showing
  that pathway features reduce confounder separability and explain when
  transfer succeeds.
\end{enumerate}

\draftnote{Replace TODO citations for Geneformer, scGPT, UCE, and
scFoundation once the bibliography is finalized. Keep v8 translational claims
out of this submission.}

\section{Benchmark Design}

GeneLab Benchmark focuses on flight-versus-ground classification under
mission shift. The full v1-v7 release spans 8 tissues: liver, gastrocnemius,
kidney, thymus, skin, eye, lung, and colon. The public OSDR source catalog
contains 24+ accessions, including excluded and track-specific rows documented
in the repository data catalog. The processed release layers contain 600+
binary/control samples.

The central evaluation unit is the mission. In leave-one-mission-out (LOMO)
tasks, one mission is held out for testing and all feature selection is
performed using training missions only. Lung and colon are single-mission
extensions and are evaluated as part of the broader v4 multi-method benchmark,
but LOMO and foundation-model summary rows must be labeled by their relevant
analysis subset.

\placeholderfigure{Benchmark schematic. The final figure should show OSDR
source data, tissue-specific task construction, mission-held-out folds,
feature surfaces, and model families.}{Figure 1 placeholder: OSDR source data
$\rightarrow$ task construction $\rightarrow$ mission-held-out evaluation.}

\subsection{Feature Surfaces and Models}

The benchmark evaluates gene-level log-normalized expression and pathway-level
features derived from Hallmark, KEGG, and combined pathway sets. The v4
multi-method surface compares PCA-LR, ElasticNet-LR, random forest, XGBoost,
SVM-linear, SVM-RBF, TabNet, and LightGBM. The foundation-model snapshot
includes Mouse-Geneformer, scGPT, UCE, and scFoundation
\citep{geneformer_todo,scgpt_todo,uce_todo,scfoundation_todo}, plus text LLM
zero-shot gene-list reasoning.

\section{Results}

\subsection{Classical Baselines Remain Strong Under Mission Shift}

Across the v4 multi-method evaluation, PCA-LR is the strongest overall
gene-level classifier, with an 8-tissue mean AUROC of 0.776; ElasticNet-LR is
second at 0.762. Table~\ref{tab:headline} gives the canonical public-facing
best-row tissue summary for v7.1.

\begin{table}[t]
\centering
\small
\caption{Canonical v7.1 v4 multi-method headline table. Significance notes
refer to best-row permutation results where available; 40 of 256 individual
evaluations are significant at $p<0.05$, and 6 of 8 tissues have at least one
significant configuration.}
\label{tab:headline}
\begin{tabular}{lclll}
\toprule
Tissue & Best AUROC & Method & Feature & Significance note \\
\midrule
Thymus & 0.948 & PCA-LR & KEGG & $p<0.05$ \\
Colon & 0.921 & PCA-LR & KEGG & $p<0.05$ \\
Lung & 0.901 & PCA-LR & Gene & $p<0.05$ \\
Kidney & 0.829 & ElasticNet-LR & Hallmark & $p<0.01$ \\
Eye & 0.823 & PCA-LR & Hallmark & not significant \\
Skin & 0.819 & PCA-LR & Gene & not significant \\
Gastrocnemius & 0.776 & PCA-LR & Gene & not significant \\
Liver & 0.670 & PCA-LR & Gene & not significant \\
\bottomrule
\end{tabular}
\end{table}

The tissue hierarchy is scientifically important. Thymus and colon lead the
best-row table, while liver is not the strongest generalizer. In the original
cross-mission transfer analysis, thymus substantially outperforms liver
(0.860 versus 0.577, permutation $p=0.001$), challenging liver-centric
assumptions about conserved spaceflight transcriptomic response.

\placeholderfigure{v4 multi-method results. The final figure should combine a
method-by-tissue heatmap with the best-row tissue summary and highlight the
PCA-LR and ElasticNet-LR baselines.}{Figure 2 placeholder: multi-method AUROC
heatmap and best-row tissue hierarchy.}

\subsection{Pathway Features Explain Tissue-Specific Transfer}

Gene-level features can strongly separate mission and hardware contexts,
whereas pathway features reduce this confounder separability. In the mission
identity task for liver, gene-level features reach macro-F1 of 1.000, while
pathway-level features fall near 0.056. This 17.8-fold contrast motivates the
central interpretation: pathway aggregation can suppress mission-specific
technical or cohort structure while retaining biological response signal.

This does not mean pathway features always outperform gene features. Instead,
they are complementary. Pathways rescue selected tissue contexts, including
kidney and eye in the original detection analysis, and improve best-row
performance in kidney, thymus, and eye in the v4 multi-method surface. Gene
features remain strongest for skin and lung.

\placeholderfigure{Gene-vs-pathway explanation. The final figure should show
confounder F1 for gene versus pathway features and selected kidney/eye pathway
rescue examples.}{Figure 3 placeholder: pathway features reduce mission
separability and rescue selected tissues.}

\subsection{Foundation Models Do Not Yet Beat the Classical Surface}

The foundation-model snapshot should be read as a benchmark-surface summary,
not a uniform 8-tissue leaderboard. Mouse-Geneformer and scGPT report
6-tissue v1 means; scFoundation and UCE rows report best single-tissue v3
results. Under these caveats, current gene-expression foundation models remain
below the classical benchmark surface: Mouse-Geneformer has a 6-tissue mean
AUROC of 0.476, scGPT has a 6-tissue mean of 0.666, and UCE and scFoundation
best single-tissue rows remain below the PCA-LR surface. Text LLMs operating on
gene-list prompts perform near chance (0.47--0.51).

Held-out validation supports the biological signal beyond the original LOMO
folds. Thymus RR-23 reaches AUROC 0.905 with $p=0.005$, and skin RR-7 reaches
AUROC 0.885 with $p<0.001$.

\placeholderfigure{Foundation model comparison and validation. The final figure
should show classical baselines versus FM/LLM rows with explicit subset labels,
plus held-out thymus and skin validation insets.}{Figure 4 placeholder:
classical baselines, FM/LLM comparison, and held-out validation.}

\section{Discussion}

GeneLab Benchmark highlights a practical gap between biological foundation
model ambition and current performance in small-sample bulk transcriptomics.
The strongest baseline is not a large pretrained model, but a carefully tuned
classical pipeline evaluated under mission-held-out splits. This does not imply
that biological foundation models are unpromising. It identifies the stress
test they must pass: robust generalization across real mission, tissue, and
batch structure.

The pathway results suggest a route forward. Generalization is not only a model
capacity problem; it is also a representation problem. Gene-level expression
captures discriminative signal but can also capture mission-specific structure.
Pathway-level aggregation can reduce confounder separability and preserve
biological response patterns that transfer across missions.

Several limitations remain. The benchmark is mouse-centered, mission folds are
small, and several foundation-model rows use different subset surfaces. The
current v7.1 submission should therefore avoid overclaiming a uniform
foundation-model leaderboard. Future work can extend the benchmark to larger
human and multi-omics settings, but translational v8 SpaceMed intervention,
countermeasure, and Mars-extrapolation claims are outside the scope of this
paper.

\paragraph{Availability.}
Source datasets are available from NASA OSDR \citep{nasa_osdr}. Code and task
definitions are available in the GeneLab Benchmark repository
\citep{genelab_benchmark_repo}. Processed feature matrices are available from
the public Hugging Face dataset card \citep{genelab_benchmark_hf}.

\paragraph{Submission status.}
This manuscript is prepared as an MLCB 2026 8-page paper. MLCB currently uses
single-blind review, asks for 11 pt font and 1 inch margins, and permits
optional PMLR publication for 8-page papers \citep{mlcb_submit_2026}. The
authors should decide whether to opt into PMLR before final submission.

\bibliographystyle{plainnat}
\bibliography{references}

\end{document}
